\section{Finite quotients of an odd-degree unitary group}
\label{sec:colors}

We now consider the first remaining family: special unitary groups defined
by division algebras of odd reduced degree.  The finite defect constructed
in Proposition~\ref{prop:finite-defect} is a quotient of a finite-index
normal subgroup of the special unitary group.  The purpose of this section
is to use a nontrivial simple quotient of that defect to construct a
homomorphism defined on the \emph{full} unitary group.  The
abelian and nonabelian alternatives will then be excluded separately.

Let $L/k$ be a quadratic extension of number fields, let $D$ be a central
simple $L$-algebra of reduced degree $n$, and let $*$ be a unitary involution
on $D$.  Write
\[
 S=\SU(D,*),\qquad U=\U(D,*),\qquad
 L^1=\{a\in L^\times:a\bar a=1\},
\]
where the bar is the nontrivial automorphism of $L/k$.  The determinant
homomorphism on $U$ is the reduced norm:
\[
 \det=\Nrd_{D/L}:U(k)\longrightarrow L^1,
 \qquad S(k)=\ker(\det).
\]
Fix $i\in L^\times$ with $\bar i=-i$.  On a dense open subset of the
$k$-vector space of Hermitian elements, the Cayley map
\begin{equation}\label{eq:colors-cayley}
 x\longmapsto (x+i)(x-i)^{-1},\qquad x=x^*,
\end{equation}
is a birational parametrization of $U$.  In particular, $U(k)$ satisfies
weak approximation: first perturb finitely many local targets into this
open chart and then approximate their inverse Cayley coordinates.

We recall the relevant local descriptions of these groups
\cite{PRbook,Reiner,KMRT}.  At a nonsplit nonarchimedean place of $L/k$,
the algebra $D$ splits.  Indeed, the existence of the unitary involution
makes its Brauer corestriction zero
\cite[Theorem~3.1(2)]{KMRT}, and corestriction for a local field
extension preserves the local Brauer invariant.  A Hermitian form of
dimension at least three over a quadratic extension of nonarchimedean
local fields is isotropic.  Thus, when $n\geq3$, the finite places where
$S$ has rank zero are split places of $L/k$.  At each such place $v$,
choose one of the two components of $D\otimes_k k_v$ and denote it by
$\mathcal D_v$.  It is a central division algebra of degree $n$ over
$k_v$, the involution exchanges the two components, and
\begin{equation}\label{eq:colors-anisotropic-factors}
 U(k_v)=\mathcal D_v^\times,
 \qquad S(k_v)=\SL_1(\mathcal D_v).
\end{equation}
Under this identification, the local determinant is the reduced norm on
$\mathcal D_v$.

For the rest of this section assume that $D$ is a division algebra and
that $n\geq3$ is odd.  Then $S$ is $k$-anisotropic.  Let $A$ be its finite
set of anisotropic nonarchimedean places and use the notation of
Proposition~\ref{prop:finite-defect}:
\begin{equation}\label{eq:colors-defect}
 R=S(k)^e,\qquad
 P_0=\overline{\delta(R)}\subset S_A,
 \qquad V_0=\delta^{-1}(P_0),\qquad V=V_0/R,
\end{equation}
where $e\geq1$, $S_A=\prod_{v\in A}S(k_v)$, and the power subgroup is
abstract.  Finally, put $A_0=\det U(k)\subset L^1$.  An empty product over
$A$ is the trivial group throughout.

\begin{proposition}[Finite quotient dichotomy]\label{prop:color-dichotomy}
Suppose that the finite group $V$ in Equation~\eqref{eq:colors-defect} is
nontrivial.  Then at least one of the following exists:
\begin{enumerate}[label=\textup{(\roman*)}]
 \item a homomorphism $\chi:U(k)\to C$ to a finite abelian group such
 that $\chi(S(k))\ne1$;
 \item a homomorphism $\phi:U(k)\to\mathcal F$ to a finite nonabelian
 simple group such that
 \[
  \phi(V_0)=\mathcal F,\qquad R\subset\ker\phi.
 \]
\end{enumerate}
\end{proposition}

The first step is an approximation statement that retains both
commutativity and exact determinants.  Its second part is also required
to retain a specified coset modulo the abstract subgroup $R$.

\subsection{Commuting approximation}

\begin{proposition}\label{prop:commuting-approximation}
With the preceding notation, the following statements hold.
\begin{enumerate}[label=\textup{(\roman*)}]
 \item Let $a,b\in A_0$.  For every $v\in A$, suppose that
 $x_v,t_v\in U(k_v)$ commute and have determinants $a_v,b_v$.
 There exist commuting $x',t\in U(k)$ with
 $\det x'=a$ and $\det t=b$ whose $A$-components approximate the
 prescribed pair arbitrarily closely.
 \item Let $\gamma\in V_0$ and $h\in U(k)$.  There exist commuting
 $x',t\in U(k)$ with
 \[
  x'\in\gamma R,\qquad \det t=\det h,
 \]
 and with $t_A$ arbitrarily close to $h_A$.
\end{enumerate}
\end{proposition}

\begin{proof}
For part~\textup{(ii)}, set $a=1$ and $b=\det h$, so that $a,b$ denote
the two prescribed determinants in either part.
We construct a regular first element $x'$ whose centralizer has enough
local determinant norms everywhere.  Lemma~\ref{lem:norm-torus} will
then produce the second element in that centralizer with its prescribed
determinant and approximations.

\paragraph{Local choices at the prescribed places.}
At $v\in A$, a commuting pair in $\mathcal D_v^\times$ lies in a common
maximal subfield.  This is because the algebra it generates is a
commutative domain in a division algebra, and hence a field, which can
be extended to a maximal subfield.  In the corresponding maximal torus
of $U(k_v)$, perturb the first element within its fixed-determinant
fiber until it is regular.  Such regular elements are dense in that
fiber: geometrically the fiber is a translate of
\[
 \{(z_1,\ldots,z_n)\in\bbG_m^n:z_1\cdots z_n=1\},
\]
and no equality $z_r=z_s$ holds identically on it.  These fibers are
smooth, so their nonempty Zariski opens are dense in the local topology.
The conjugation map around a regular semisimple element is a local
submersion.  Consequently, sufficiently small changes of this first
element have centralizer tori obtained by conjugation with elements
arbitrarily close to $1$.  Conjugating the second element with the same
elements retains its determinant and its prescribed approximation.

For part~\textup{(ii)}, take the first local element sufficiently close
to $1$ in a maximal torus containing $h_v$, and with determinant $1$.
The choice near $1$ is compatible with the closure of $\gamma R$ because
$\gamma\in V_0$.  The same regular perturbation applies.  Hence in both
parts we have nonempty open conditions for the first element such that
its local centralizer admits the required second element.

Enlarge the finite set of prescribed places to include all archimedean
places, places of bad reduction or small residue characteristic, places
ramified in $L/k$, and places where $a$ or $b$ is not a unit.  At every
additional place choose a maximal torus whose determinant image contains
both required determinants.  At nonsplit places the algebra splits, and
a diagonal unitary torus has determinant image equal to the local
norm-one group.  At a split finite place outside $A$, write the chosen
component of the algebra as $M_l(\Delta)$.  If $\Delta$ is nontrivial,
then $l\geq2$: use an unramified maximal subfield of $\Delta$ in one
diagonal block and a totally ramified maximal subfield in another.
Their product norm image contains all units and an element of valuation
one, and is therefore all of $k_v^\times$.  Complete these two blocks
to a maximal \'etale subalgebra.  If $\Delta=k_v$, the usual diagonal
torus suffices.  At split archimedean places the algebra is split because
its degree is odd.  In each case we may choose the first element regular
with the specified determinant and then take a sufficiently small
opening around it.

We also prescribe finitely many auxiliary split places, outside all
exceptions, at which the algebra splits and $a,b$ are units.  Choose
unramified maximal tori having every possible permutation cycle type in
$S_n$, one type at each such place.  Norms from their unramified field
factors cover the unit group.  These conditions will force the joint
Galois group needed by Lemma~\ref{lem:norm-torus}.

In part~\textup{(i)}, the determinant-$a$ slice is a rational translate
of $S$ because $a\in A_0$.  Weak approximation for $S$, followed by
Proposition~\ref{prop:power-approximation} on that translate, permits all
the local openings just chosen.  In part~\textup{(ii)}, apply that
proposition in the coset $\gamma R$.  The openings near $1$ at $A$ are
allowed by Equation~\eqref{eq:colors-defect} and the closure statement
in Proposition~\ref{prop:finite-defect}; the other local openings are
unrestricted.  We now check the reduction and pole conditions required
by Proposition~\ref{prop:power-approximation}.

\paragraph{Integral reductions and single poles.}
At an integral reduction outside the prescribed set, exclude matrices
whose characteristic polynomial has no simple root over an algebraic
closure of the residue field.  This exclusion has geometric codimension
at least two in a fixed nonzero determinant slice of $\GL_n$.
Indeed, if every eigenvalue has multiplicity at least two, there are at
most $\lfloor n/2\rfloor$ distinct eigenvalues.  After fixing the
determinant, they contribute at most $\lfloor n/2\rfloor-1$ parameters.
For each Jordan type with fixed eigenvalues, the orbit dimension is at
most $n^2-n$.  Thus the excluded set has dimension at most
\[
 n^2-n+\lfloor n/2\rfloor-1
 \leq (n^2-1)-2.
\]
There are only finitely many Jordan types, so the same bound holds for
their union and its closure.  These conditions descend to the unitary
determinant slice and spread out after finitely many model exceptions.

A simple absolute residue root belongs to a residue-field irreducible
factor of multiplicity one.  Hensel lifting gives an unramified field
factor of the algebra generated by the first element.  We do not need
the simple root itself to lie in the ground residue field.

At a single pole in Proposition~\ref{prop:power-approximation}, the place
splits in the fixed splitting field and the first element has the form
\[
 x'=g_L\,\diag(t^{e d_1},\ldots,t^{e d_n})\,g_R,
 \qquad d_1<\cdots<d_n,
\]
where $\ord_v(t)>0$ and the two side factors are integral and invertible.
Require every principal minor on the first $j$ indices of $g_Rg_L$ to
be nonzero modulo the place, for $1\leq j\leq n$.  This is a nonempty
open condition on each pole hyperplane: an undamaged basic parameter
list in Proposition~\ref{prop:power-approximation} is dominant, so its
varying side factor can meet the finitely many principal-minor opens.
If $c_j$ is the $j$th elementary characteristic coefficient, principal
minor expansion gives
\begin{equation}\label{eq:colors-pole-valuations}
 \ord_v(c_j)=e(d_1+\cdots+d_j)\ord_v(t).
\end{equation}
The term from the first $j$ indices is the unique term of least
valuation.  The successive Newton slopes are strictly distinct and all
segments have horizontal length one.  The characteristic polynomial
therefore splits into distinct linear factors over $k_v$.

The additional denominator places in the approximation procedure admit
the same determinant-compatible local tori described above.  At the
later finite model exceptions, integral unit parameters giving regular
reduction provide the required openings.  This checks all hypotheses
of Proposition~\ref{prop:power-approximation}.

\paragraph{The global centralizer and its norm torsor.}
Choose $x'$ by that proposition, also imposing global regularity.  Since
$D$ is division,
\[
 E=L(x')\subset D
\]
is a maximal field of degree $n$ over $L$.  It is stable under $*$ because
$(x')^*=(x')^{-1}$.  Put $F=\{z\in E:z^*=z\}$, the fixed field.
Then $E=LF$ and
$[F:k]=n$.

Let $\Gamma\subset S_n\times C_2$ be the Galois group of the normal
closure of $F$ together with $L$, acting on the embeddings of $F$ and on
$L$.  At the auxiliary places split in $L$, the prescribed cycle types
show that the kernel of $\Gamma\to C_2$ meets every conjugacy class of
$S_n$.  A proper subgroup of a finite group cannot have this property:
the transitive action on its cosets would have a derangement, by
averaging the number of fixed cosets, whereas every conjugacy class
meeting the subgroup would rule out all derangements.  Hence that
kernel is $S_n\times\{1\}$.  The projection to $C_2$ is surjective, so
\[
 \Gamma=S_n\times C_2.
\]

The unitary torus in the centralizer of $x'$ is
\[
 T_E=\Res_{F/k}\bigl(\Res^1_{E/F}\bbG_m\bigr),
\]
and its determinant is $N_{E/L}:T_E\to\Res^1_{L/k}\bbG_m$.
At all prescribed places its image contains $b_v$ by construction.
At every other integral place, take the full $*$-orbit of the unramified
field factor of $E$ found above; all its fields are unramified over $k_v$
since $L/k$ is unramified here.  Its fixed algebra is an unramified field
factor $F_j/k_v$ of $F\otimes_k k_v$: it is the fixed field when the
factor is stable, and the diagonal fixed algebra when two factors are
exchanged.  The corresponding $*$-stable algebra is
$F_j\otimes_{k_v}(L\otimes_k k_v)$, possibly split, so no individual
$E$-factor is required to be $*$-stable.  Restrict the norm
map to this factor.  It is a surjective map of unramified tori with
connected kernel.  To check connectedness, over an algebraic closure
the dual map sends a generator to a vector with entries $1$ or $-1$;
this vector is primitive.  Lang's theorem on the residue-field tori and
smooth lifting show that the norm map is onto integral points.  As $b_v$
is a unit, it is a local determinant norm.  At a single pole the torus
splits by Equation~\eqref{eq:colors-pole-valuations}, and the same
conclusion holds without an integrality restriction.

Thus the fiber $N_{E/L}^{-1}(b)$ has points everywhere locally.  Its
acting kernel is exactly the torus of Lemma~\ref{lem:norm-torus}, with
the joint Galois group just verified.  That lemma gives a rational point
$t$ in the fiber approximating the chosen local second elements at $A$.
It belongs to the centralizer of $x'$, so $x'$ and $t$ commute.  The
construction retained the required determinant and, in part~\textup{(ii)},
the coset $\gamma R$.  This proves both assertions.
\end{proof}

\subsection{A simple quotient preserved by the full unitary group}

The passage from class preservation to a quotient of the ambient group
follows the inner-type argument of Rapinchuk--Potapchik
\cite[Corollary~2.3 and Theorem~2.4]{RapinchukPotapchik1996}.
Here Proposition~\ref{prop:commuting-approximation} supplies the required
unitary input. Conjugation by $U(k)$ preserves $R$, since $R$ is
characteristic in $S(k)$.  It also preserves $P_0$ and $V_0$ by passage to local closures.
In fact the resulting action on $V$ preserves each conjugacy class.
To see this, take $\gamma\in V_0$ and $h\in U(k)$.  By
Proposition~\ref{prop:commuting-approximation}\textup{(ii)}, choose
$x'\in\gamma R$ commuting with $t$, with $\det t=\det h$ and $t_A$
sufficiently close to $h_A$ that $t h^{-1}\in V_0$.  In $V$, conjugation
by $t$ fixes the image of $x'$, whereas $t h^{-1}$ induces an inner
automorphism.  Conjugation by $h$ consequently carries $\gamma R$ to a
conjugate of itself.

If $V\ne1$, choose a maximal proper normal subgroup of $V$ and denote
the resulting simple quotient by $B$.  Its kernel is invariant under
the $U(k)$ action because a normal subgroup is a union of conjugacy
classes.  The induced automorphisms of $B$ are class-preserving.  If $B$
is nonabelian, the theorem of Feit and Seitz
\cite[Theorem~C]{FeitSeitz} says that every such automorphism is inner.
Since a nonabelian simple group has trivial center, the identification
$\Inn(B)=B$ gives
\[
 \phi:U(k)\longrightarrow B.
\]
Its restriction to $V_0$ is the quotient map to $B$, and $R$ acts
trivially.  This is alternative~\textup{(ii)} of
Proposition~\ref{prop:color-dichotomy}.

It remains to treat the case that $B$ is cyclic of prime order.  Let
$R'$ be its kernel upstairs, so that
\begin{equation}\label{eq:colors-central-quotient}
 R\subset R'\subset V_0,\qquad B=V_0/R'.
\end{equation}
Class preservation on this abelian quotient means that the action is
trivial.  Thus $R'$ is normal in $U(k)$ and $B$ is central in $U(k)/R'$.
The rest of the section uses a nonzero character of this central
subgroup to construct a finite character of $U(k)$ that is nonzero on
$S(k)$.

\subsection{Local abelianization and the metaplectic input}

We need two local facts: the continuous characters of the compact groups
$S(k_v)$, and a splitting theorem for their finite product.  The first
fact also identifies the closed subgroup generated by commutators of
the full local multiplicative group.

\begin{lemma}\label{lem:colors-local-abelianization}
Let $F_v$ be a nonarchimedean local field of characteristic zero, with
residue field $\bbF_q$, and let $\mathcal D$ be a central division algebra
over $F_v$ of odd degree $n\geq3$.  Put $S_v=\SL_1(\mathcal D)$.
Reduction induces an isomorphism
\[
 S_v/\overline{[S_v,S_v]}
 \simeq T_v:=\ker\bigl(N_{\bbF_{q^n}/\bbF_q}:
                       \bbF_{q^n}^\times\to\bbF_q^\times\bigr).
\]
Moreover,
\[
 \overline{[\mathcal D^\times,\mathcal D^\times]}=S_v.
\]
If an element of $\mathcal D^\times$ has reduced-norm valuation $r$,
its action on $T_v$ is $\sigma^r$, where $\sigma:z\mapsto z^{q^h}$
for some $h$ relatively prime to $n$.
\end{lemma}

\begin{proof}
Use the cyclic description of a division algebra over a local field
\cite{Reiner}.  There is an unramified maximal subfield $E_v/F_v$ and
a prime element $\Pi$ such that conjugation by $\Pi$ induces a generator
$\sigma$ of $\Gal(E_v/F_v)$.  The residue algebra is $\bbF_{q^n}$ and
$\sigma$ acts there by $q^h$ with $(h,n)=1$.  Normalize valuations so
that $\ord_v(\Nrd\Pi)=1$.  Every element of $S_v$ is integral and its
residue lies in $T_v$.  The Teichm\"uller representatives in $E_v$ lift
$T_v$ to a subgroup of $S_v$, so reduction onto $T_v$ is surjective.

Let $\mathfrak P=(\Pi)$ be the maximal ideal of the maximal order of
$\mathcal D$, and set
\[
 S_r=S_v\cap(1+\mathfrak P^r),\qquad r\geq1.
\]
In leading-coefficient coordinates, the quotients of this filtration
are
\begin{equation}\label{eq:colors-filtration}
 S_r/S_{r+1}\simeq
 \begin{cases}
  (\bbF_{q^n},+),&n\nmid r,\\
  \ker\Tr_{\bbF_{q^n}/\bbF_q},&n\mid r.
 \end{cases}
\end{equation}
For completeness, the norm filtration is
\[
 \Nrd(1+\mathfrak P^r)=1+\mathfrak p^{\lceil r/n\rceil}.
\]
The inclusion follows from the valuations of the characteristic
coefficients; the reverse inclusion follows already from norms of
principal units in the unramified maximal subfield.  If $n\nmid r$,
the norm images for $r$ and $r+1$ agree, so any leading coefficient can
be corrected at the next level to have norm one.  If $n\mid r$, the
first norm coefficient is the residue trace, and the next-level norm
image gives exactly the trace-zero restriction in
Equation~\eqref{eq:colors-filtration}.  The finite-field trace is onto
also when the residue characteristic divides $n$.

For $n\nmid r$, commute an element at level $r$ with a Teichm\"uller
element of $T_v$ not fixed by $\sigma^r$.  Such an element exists:
$T_v$ is cyclic of order $(q^n-1)/(q-1)$ and is not contained in the
multiplicative group of a proper subfield.  On the leading coefficient
this commutator is multiplication by a nonzero residue scalar.
Thus commutators fill $S_r/S_{r+1}$ in this case.

For $n\mid r$, use brackets between levels $1$ and $r-1$.  Both levels
have unrestricted leading coefficients because $n\geq3$.  Taking the
coefficient $1$ in degree $1$ makes the bracket coefficients contain
\[
 (\sigma-1)\bbF_{q^n}=\ker\Tr_{\bbF_{q^n}/\bbF_q}.
\]
Hence commutators again fill the quotient.  Successive approximation
now puts all of $S_1$ in $\overline{[S_v,S_v]}$.  The reverse inclusion
holds because the residue quotient $T_v$ is abelian.  This proves the
first assertion.

Finally, commutators of $\Pi$ with unramified units have residues
$\sigma(z)/z$, which fill $T_v$ by finite-field Hilbert 90.  Together
with the already obtained subgroup $S_1$, they generate a dense subgroup
of $S_v$.  Every commutator in $\mathcal D^\times$ has reduced norm one,
proving the second assertion.  A local unit acts trivially on the
commutative residue field, while $\Pi$ acts by $\sigma$.  Writing an
arbitrary element as a unit times $\Pi^r$ gives the last statement.
\end{proof}

Here is the precise metaplectic consequence we use.  Let
$I=\bbR/\bbZ$, written additively.  If $G$ is a simply connected
absolutely almost simple group over a number field, not of type $A_1$,
and $A'$ is a finite set of nonarchimedean places where $G$ is anisotropic,
then the restriction map
\begin{equation}\label{eq:colors-metaplectic}
 H^2_{\mathrm m}\!\left(\prod_{v\in A'}G(k_v),I\right)
 \longrightarrow H^2\bigl(G(k),I\bigr)
 \quad\text{is injective}.
\end{equation}
This is the specialization $V_2=\varnothing$ of
\cite[Proposition~2.6, p.~114]{PRmetaplectic}. If $A'$ is nonempty,
the local classification forces the absolute type to be $A$, so the
type-$D$ exclusion inherited from Proposition~2.4 there is harmless.
If $A'$ is empty, the source is $H^2_{\mathrm m}(1,I)=0$.
The source uses measurable
cohomology for locally compact groups and abstract cohomology on the
rational group.  In degree two these classify, respectively,
topological central circle extensions with Borel sections and abstract
central extensions; see \cite[Section~3.2]{PRsurvey}.  Thus a topological
central circle extension of the displayed finite local product that
splits over the abstract rational group has a continuous homomorphic
section.  Indeed, the zero measurable class gives a Borel homomorphic
section, and Borel homomorphisms between these locally compact
second-countable groups are continuous.

We apply only Equation~\eqref{eq:colors-metaplectic} with $G=S$ and
$A'=A$.  Here $n\geq3$ excludes type $A_1$, and every support place is
nonarchimedean and anisotropic.  In particular, this invocation involves
neither a comparison with real-place metaplectic kernels nor an
identification of a congruence kernel requiring the normal-subgroup
assertion we are proving.

\subsection{A circle extension of the full local determinant group}

Give $A_0$ the discrete topology and define
\begin{equation}\label{eq:colors-Q}
 Q=\{(y,a)\in U_A\times A_0:\det y=a_A\},
 \qquad U_A=\prod_{v\in A}U(k_v).
\end{equation}
Its determinant kernel is $S_A$.  The group $Q$ is locally compact and
second countable: each of its countably many determinant fibers is a
translate of the compact group $S_A$.  There is a natural homomorphism
\[
 \rho:U(k)\longrightarrow Q,\qquad u\longmapsto(u_A,\det u).
\]
The subgroup $P_0\subset S_A$ is open and normal in $Q$.  Normality
follows first under the dense group $\rho(U(k))$ and then under $Q$.
In fact $\rho(U(k))$ is dense in every determinant fiber by weak
approximation for $S$.  The same argument, applied modulo the open
subgroup $P_0$, gives an isomorphism of discrete groups
\begin{equation}\label{eq:colors-discrete-quotient}
 Q/P_0\simeq U(k)/V_0.
\end{equation}

The central quotient in Equation~\eqref{eq:colors-central-quotient}
provides the discrete central extension
\[
 1\longrightarrow B\longrightarrow U(k)/R'
 \longrightarrow U(k)/V_0\longrightarrow1.
\]
Pull it back along Equation~\eqref{eq:colors-discrete-quotient} and push
it out by an injective character $\iota:B\hookrightarrow I$.  We obtain
a locally compact second-countable central extension
\begin{equation}\label{eq:colors-circle-extension}
 1\longrightarrow I\longrightarrow\widetilde Q
 \xrightarrow{p}Q\longrightarrow1.
\end{equation}
The pullback description gives continuous local sections, and it also
gives a specified abstract lift
\[
 \ell:U(k)\longrightarrow\widetilde Q,
 \qquad p\circ\ell=\rho.
\]
Over $P_0$ there is a canonical continuous homomorphic section
$c:P_0\to\widetilde Q$, obtained by taking the identity element in the
discrete quotient $U(k)/R'$.  For $u\in V_0$, these descriptions give
\begin{equation}\label{eq:colors-canonical-lift}
 \ell(u)=c(u_A)\,\iota(uR'),
\end{equation}
where the last factor denotes the central element of $\widetilde Q$
corresponding to $\iota(uR')\in I$.

\begin{lemma}\label{lem:colors-commuting-lifts}
If two elements of $Q$ commute, any lifts of them to $\widetilde Q$
commute.
\end{lemma}

\begin{proof}
The two determinant coordinates are elements $a,b\in A_0$.  By
Proposition~\ref{prop:commuting-approximation}\textup{(i)}, their local
components are limits of commuting rational pairs with these exact
determinants.  The specified lifts under $\ell$ of each rational pair
commute.  For a commuting pair in the base of a central extension, the
commutator of lifts is independent of the choices of lifts.  Use the
continuous local sections of Equation~\eqref{eq:colors-circle-extension}
to pass these commutators to the limit.  Their limit is $1$, as required.
\end{proof}

Restriction of Equation~\eqref{eq:colors-circle-extension} to $S_A$
splits abstractly over $S(k)$ by $\ell$.  Equation~\eqref{eq:colors-metaplectic}
therefore gives a continuous homomorphic section
\[
 j:S_A\longrightarrow\widetilde Q.
\]
Our remaining task is to modify $j$ so that this splitting extends from
$S_A$ to all of $Q$.  We first make $j(S_A)$ normal and then eliminate
the commutator pairing on the determinant quotient $A_0$.

\subsection{Making the splitting equivariant}

For $q\in Q$, choose a lift $\widetilde q$ and define the discrepancy
character $d_q:S_A\to I$ by
\begin{equation}\label{eq:colors-discrepancy}
 \widetilde q\,j(s)\,\widetilde q^{-1}
   =j(qsq^{-1})\,d_q(s).
\end{equation}
We again regard values in $I$ as central factors.  This is a continuous
character of $s$, independent of the lift of $q$.  It is zero for
$q\in S_A$, since there a lift differs from $j(q)$ by a central element.
Composition of conjugations shows that the $d_q$ form a cocycle for the
conjugation action.  More explicitly,
\[
 d_{q_1q_2}(s)=d_{q_1}(q_2sq_2^{-1})+d_{q_2}(s).
\]
The cocycle and its action factor through $Q/S_A=A_0$.

By Lemma~\ref{lem:colors-local-abelianization}, its coefficient group is
the finite group
\[
 \Hom_{\mathrm{cont}}(S_A,I)
   =\prod_{v\in A}\Hom(T_v,I).
\]
On the $v$-component the action of $a\in A_0$ is the power of $\sigma_v$
specified by $\ord_v(a_v)$.  The corresponding component of the
discrepancy is zero whenever $\ord_v(a_v)=0$.  Indeed, choose a
representative of $a$ in $Q$ whose $v$-component is an unramified unit
of reduced norm $a_v$.  Such a unit exists by surjectivity of unramified
norms on units.  It commutes with every Teichm\"uller lift of $T_v$,
viewed as an element of $S_A$ supported at $v$.  Lemma~\ref{lem:colors-commuting-lifts}
then makes Equation~\eqref{eq:colors-discrepancy} zero on these lifts,
which determine the character.

It follows that both this component of the cocycle and its action factor
through the single valuation map $A_0\to\bbZ$.  This map is onto:
weak approximation in $U(k)$ allows its $v$-component to be close to
a local element of reduced-norm valuation one.  Let $d_v$ be the
discrepancy for such a generator.  It vanishes on $T_v^{\sigma_v}$.
Indeed, the Teichm\"uller lifts of these fixed residue elements are
central in $\mathcal D_v^\times$, so the same commuting-lifts test
applies for any representative of the generator.

For a finite abelian group $T$ with automorphism $\sigma$, the image of
\[
 \sigma^*-1:\Hom(T,I)\longrightarrow\Hom(T,I)
\]
is exactly the characters annihilating $T^\sigma$.  This follows by
duality from the exact sequence for $\sigma-1$ on $T$, or directly by
extending a character from its image to $I$.  Consequently $d_v$ is a
coboundary.  Choose a correcting character in each component and
replace $j(s)$ by $j(s)\lambda(s)$ for the resulting continuous character
$\lambda:S_A\to I$.  Equation~\eqref{eq:colors-discrepancy} then has
zero discrepancy for every $q$.  Thus we may, and do, assume that
\begin{equation}\label{eq:colors-equivariant-section}
 \widetilde q\,j(s)\,\widetilde q^{-1}=j(qsq^{-1})
 \quad(q\in Q,\ s\in S_A),
\end{equation}
so that $j(S_A)$ is normal in $\widetilde Q$.

\subsection{Removing the determinant commutator pairing}

The quotient by this normal subgroup is a central extension
\begin{equation}\label{eq:colors-determinant-extension}
 1\longrightarrow I\longrightarrow\widetilde Q/j(S_A)
 \longrightarrow A_0\longrightarrow1.
\end{equation}
Since $A_0$ is abelian, commutators in this extension give an alternating
biadditive pairing, written multiplicatively in its arguments,
\[
 \omega:A_0\times A_0\longrightarrow I.
\]
By Lemma~\ref{lem:colors-commuting-lifts}, $\omega(a,b)=0$ whenever,
at every $v\in A$, the two determinants $a_v,b_v$ are norms of elements
of one common commutative subfield of $\mathcal D_v$.

In particular, if every $a_v$ is an $n$th power, represent it by a
central scalar at each place; it commutes with any representative of
$b_v$.  Thus $\omega$ factors through the finite group
\begin{equation}\label{eq:colors-power-quotient}
 K_A=\prod_{v\in A}k_v^\times/(k_v^\times)^n.
\end{equation}
The map $A_0\to K_A$ is onto.  Local reduced norms are surjective, and
weak approximation in $U(k)$ realizes any prescribed classes because
the local power subgroups are open.  Mixed terms between two distinct
factors of Equation~\eqref{eq:colors-power-quotient} vanish: use central
representatives for each locally trivial power class.  We may therefore
treat the pairing one place at a time.

Fix $v$, suppress it from the notation, and let $\pi$ be a uniformizer.
Unit-unit terms vanish because units are norms from the unramified
maximal subfield.  Alternation kills uniformizer-uniformizer terms, so
the local pairing is specified by the character
\[
 u\longmapsto\omega_v(\pi,u),\qquad u\in\mathcal O_v^\times.
\]
This character vanishes on $1+\mathfrak p_v$.  To prove it first take
$c\in k_v$ of valuation one.  The Eisenstein extension defined by
$\theta^n=c$ has degree $n$, is totally ramified, and embeds in
$\mathcal D_v$ by the local invariant criterion.  Since $n$ is odd,
\[
 N(\theta)=c,\qquad N(1+\theta)=1+c.
\]
The common-subfield test gives $\omega_v(c,1+c)=0$; writing $c$ as a
uniformizer times a unit and using unit-unit vanishing gives
$\omega_v(\pi,1+c)=0$.  If $\ord_v(d)\geq2$, choose any such $c$ and
write
\[
 1+d=\frac{(1+c)(1+d)}{1+c}.
\]
Both numerator and denominator have the form $1+c'$ with
$\ord_v(c')=1$.  This proves vanishing on all principal units.
The remaining parameter is therefore a character of
$\bbF_q^\times$ killed by $n$.

We still have freedom to change $j$ by characters of $S_A$ invariant
under $Q$.  At the fixed place these are the $\sigma$-invariant
characters of $T_v$.  Such a change preserves
Equation~\eqref{eq:colors-equivariant-section}.  Its effect on the
pairing is evaluation, up to the fixed sign convention for commutators,
on the residue of the ordinary local commutator of determinant
representatives.  For a prime element $\Pi$ and an unramified unit with
residue $r$, that residue is
\[
 \frac{\sigma(r)}{r}.
\]
The residue norm of $r$ is the residue of the unit determinant $u$.
Replacing $r$ by another norm preimage changes this expression by an
element of $(\sigma-1)T_v$.  Hence we have a well-defined map
\begin{equation}\label{eq:colors-residue-commutator}
 \bbF_q^\times\longrightarrow T_v/(\sigma-1)T_v,
 \qquad N(r)\longmapsto\sigma(r)/r.
\end{equation}

This map is onto.  To see this explicitly, put
\[
 M=\frac{q^n-1}{q-1},\qquad d=\gcd(q-1,n),
\]
and choose a generator $g$ of $\bbF_{q^n}^\times$.  The group $T_v$ is
generated by $g^{q-1}$.  Since $(h,n)=1$,
\[
 |T_v/(\sigma-1)T_v|
 =\gcd(M,q^h-1)=\gcd(M,q-1)=d.
\]
The image of the generator $g^M$ of $\bbF_q^\times$ in
Equation~\eqref{eq:colors-residue-commutator} is represented, in the
generator $g^{q-1}$, by the exponent
\[
 \frac{q^h-1}{q-1}\equiv h\pmod d.
\]
It is a generator because $(h,d)=1$.  Dualizing this surjection shows
that invariant characters of $T_v$ produce exactly all residue unit
characters killed by $n$.  They can therefore cancel the remaining
parameter $u\mapsto\omega_v(\pi,u)$.  A unit adjustment between
$\Nrd(\Pi)$ and the chosen $\pi$ makes no difference, since unit-unit
terms already vanish.

Perform this adjustment at every place of $A$.  The new section still
has normal image, and the commuting-lifts test still holds; in
particular, all the preceding reductions of the pairing remain valid.
Its only possible parameters have now been cancelled, so $\omega=0$.
The extension in Equation~\eqref{eq:colors-determinant-extension} is
therefore abelian.  The circle group is divisible and hence injective
as an abelian group, so this extension splits.  A homomorphism from the
discrete group $A_0$ is automatically continuous.

Let $C_0$ be the inverse image in $\widetilde Q$ of the image of such
a section of Equation~\eqref{eq:colors-determinant-extension}.  Then
$C_0\cap I=1$ and $p(C_0)=Q$: for any lift of an element of $Q$, its
central coordinate in the quotient can be removed using $I$.
Consequently $p:C_0\to Q$ is a group isomorphism.  Locally on each
determinant fiber its inverse is a translate of $j$, so the inverse is
continuous.  We have proved that Equation~\eqref{eq:colors-circle-extension}
has a continuous homomorphic section
\begin{equation}\label{eq:colors-global-section}
 s:Q\longrightarrow\widetilde Q.
\end{equation}

\subsection{Extracting a nonzero finite character}

Compare the continuous section in Equation~\eqref{eq:colors-global-section}
with the specified abstract lift of rational points.  Their difference
is a character
\begin{equation}\label{eq:colors-difference-character}
 \chi_0:U(k)\longrightarrow I,
 \qquad \ell(u)=s(\rho(u))\,\chi_0(u).
\end{equation}
It has finite image on $S(k)$.  Indeed, for $u\in R'$ the canonical
formula in Equation~\eqref{eq:colors-canonical-lift} reads
$\ell(u)=c(u_A)$.  The difference of $c$ and $s|_{P_0}$ is a continuous
character of the profinite group $P_0$.  Its image in the circle is
finite: a quotient of a profinite group is profinite, whereas a closed
subgroup of the circle is either finite or the whole circle.  Thus
$\chi_0(R')$ is finite.  Since $R'$ has finite index in $S(k)$, the image
$\chi_0(S(k))$ is a finite union of cosets of that finite subgroup and
is finite.

It is also nonzero.  Choose $\beta\ne1$ in $B$ and $u\in V_0$
representing it.  Because $R$ is dense in $P_0$ and $u_A\in P_0$,
there is a sequence $r_j\in R$ such that
\[
 (ur_j)_A\longrightarrow1.
\]
All $ur_j$ still represent $\beta$ modulo $R'$.  If $\chi_0$ vanished
on $S(k)$, Equations~\eqref{eq:colors-canonical-lift} and
\eqref{eq:colors-difference-character} would give
\[
 c((ur_j)_A)\,\iota(\beta)=s(\rho(ur_j)).
\]
Both continuous sections tend to the identity, which would imply
$\iota(\beta)=0$, contrary to injectivity of $\iota$.

Finally we make the character finite on all of $U(k)$ without changing
its restriction to $S(k)$.  The multiplicative group of a number field
is free abelian modulo its finite roots of unity: use the ideal map and
the finite generation of the unit group.  Therefore its subgroup $A_0$
has the form
\[
 A_0=T\oplus\bigoplus_{j\in J}\bbZ a_j,
\]
where $T$ is finite and $J$ may be infinite.  Choose $u_j\in U(k)$
with $\det u_j=a_j$, and choose a character $\psi:A_0\to I$ satisfying
$\psi(a_j)=\chi_0(u_j)$ and $\psi|_T=0$.  Then
\[
 \chi=\chi_0-\psi\circ\det
\]
vanishes on all the chosen $u_j$ and equals $\chi_0$ on $S(k)$.
If $m$ kills the finite group $\chi_0(S(k))$ and $t$ kills $T$, then
$mt$ kills $\chi(U(k))$: remove the free determinant part using the
$u_j$, and the $t$th power of the remaining element lies in $S(k)$.
The subgroup of the circle killed by $mt$ is finite.  Thus $\chi$ is a
finite character and remains nonzero on $S(k)$.

When $A=\varnothing$, the same construction has $Q=A_0$ and trivial
$S_A$.  Commuting approximation makes the circle extension abelian
immediately, and all subsequent arguments apply unchanged.  We have
proved alternative~\textup{(i)} in the remaining case, completing the
proof of Proposition~\ref{prop:color-dichotomy}.

The finite defect can therefore be nontrivial only if the full unitary
group admits one of the two homomorphisms in that proposition.  The
character argument in Section~\ref{sec:characters} and the color argument
in Sections~\ref{sec:palettes}--\ref{sec:finite-groups} exclude these
alternatives, retaining the division and odd-degree hypotheses fixed here.
